\begin{trapbox}
\textbf{陷阱 1｜没有规定 $A,B$ 的方向，就直接背 $AF/BF$。}
同一条直线交换端点，长度比会取倒数。
\textbf{纠正：}先写 $A=F+ru$、$B=F-su$，其中 $r,s>0$。
\par\smallskip\noindent
\textbf{陷阱 2｜把倾斜角看成与 $x$ 轴夹的锐角。}
倾斜角 $\theta$ 的取值范围是 $[0,\pi)$。
$k<0$ 时为钝角，必须有 $\cos\theta<0$；
不能用 $+1/\sqrt{1+k^2}$ 代入所有斜率。
\par\smallskip\noindent
\textbf{陷阱 3｜斜率为零、斜率不存在时强行使用 $y=k(x-c)$。}
$k=0$：$\theta=0$，$AF=a-c$、$BF=a+c$；
竖直线无斜率：$\theta=\pi/2$，$AF=BF=b^2/a$。
\textbf{纠正：}统一使用方向向量 $(\cos\theta,\sin\theta)$，不丢竖直情形。
\par\smallskip\noindent
\textbf{陷阱 4｜把 $b^2$ 写成 $a^2$ 或把 $e$ 写成 $c$。}
长度分式的分子是 $b^2=a^2-c^2$，
分母是 $a\pm c\cos\theta$；
比例分式的分子分母各除以 $a$ 后，才出现 $e=c/a$。
\par\smallskip\noindent
\textbf{陷阱 5｜``反方向''漏掉负号或把负根当成负距离。}
$B=F-su$，不是 $F+su$；参数法的 $t_-<0$
只是有向参数，真正长度 $BF=-t_->0$。
\par\smallskip\noindent
\textbf{陷阱 6｜右焦点与左焦点套同一套符号。}
沿相同正向极轴规定端点时，右焦点
$AF=b^2/(a+c\cos\theta)$；左焦点
$AF_1=b^2/(a-c\cos\theta)$。
可用水平长轴特例验证：右焦点朝右仅 $a-c$，
左焦点朝右却是 $a+c$。
\par\smallskip\noindent
\textbf{陷阱 7｜使用第二定义时，未确认准线在哪一侧。}
椭圆右准线 $x=a^2/c>a$，始终在椭圆右边，
所以点 $P$ 到准线的距离是 $a^2/c-x_P$，而非相反。
左准线则是 $x=-a^2/c$，方向不同。
\par\smallskip\noindent
\textbf{陷阱 8｜把 $AF+BF$ 写成 $|AF-BF|$。}
$A,F,B$ 按图在同一条直线上，且焦点 $F$ 在两端点之间：
$AB=AF+BF$。\textbf{别把椭圆第一定义 $AF_1+AF=2a$}
误当成焦点弦的两段之和恒为 $2a$（一般不是）。
\par\smallskip\noindent
\textbf{陷阱 9｜逆向给比值时，忘记检查是否可实现。}
若给定 $R=AF/BF>0$，则
$\cos\theta=(1-R)/[e(1+R)]$。
必须满足 $-1<\cos\theta\le1$（由于 $\theta\in[0,\pi)$）。
若题目只说无序比例，还要考虑端点交换与对称的直线。
\end{trapbox}
